% Part: second-order-logic
% Chapter: syntax-and-semantics

\documentclass[../../../include/open-logic-chapter]{subfiles}

\begin{document}

\olchapter{sol}{syn}{વાક્યરચના અને અર્થવિજ્ઞાન}

\begin{editorial}
હાલમાં આ પ્રકરણમાં દ્વિતીય-ક્રમ તર્કશાસ્ત્રની મૂળભૂત વાક્યરચના અને
અર્થવિજ્ઞાન આવરી લેવાયાં છે. પ્રકરણ તરીકે તે હજી ઘણું ટૂંકું છે.
દ્વિતીય-ક્રમ તર્કશાસ્ત્રની !!{derivation} પદ્ધતિઓની ચર્ચા કરી શકાય તે
માટે દ્વિતીય-ક્રમ ચલોનું પ્રતિસ્થાપન પણ સમજાવવું પડશે; તેમાં કેટલીક
સૂક્ષ્મ મુશ્કેલીઓ છે.
\end{editorial}

\olimport{introduction}

\olimport{terms-formulas}

\olimport{satisfaction}

\olimport{semantic-notions}

\olimport{expressive-power}

\olimport{inf-count}

\OLEndChapterHook

\end{document}
